\documentclass[11pt,a4paper]{article}
\usepackage[margin=0.85in,top=1in,bottom=1in]{geometry}
\usepackage{amsmath,amssymb,amsfonts}
\usepackage{graphicx}
\usepackage{float}
\usepackage{enumitem}
\usepackage{microtype}
\usepackage{caption}
\usepackage[hidelinks]{hyperref}
\usepackage[most,breakable]{tcolorbox}
\tcbuselibrary{breakable,skins}
\usepackage[utf8]{inputenc}
\usepackage[T1]{fontenc}
\usepackage{lmodern}
\captionsetup{font=small,labelfont=bf,skip=4pt}
\setlist{leftmargin=1.5em,topsep=2pt}

%% Breakable card — prevents content cutoff across pages
\newtcolorbox{solcard}[3]{
  breakable, enhanced,
  colback=white, colframe=gray!50, arc=2pt,
  title={\textbf{#1}\hfill\textsf{\small #3\,pts}},
  fonttitle=\normalsize\sffamily\bfseries,
  before upper={\small\textit{#2}\par\smallskip\hrule height 0.4pt\smallskip},
  top=4pt, bottom=6pt, left=7pt, right=7pt,
  parbox=false,
}

\title{\textbf{Magnetic trap --- Solution}}
\author{\normalsize X23 (2023) — English translation}
\date{}
\begin{document}\maketitle\vspace{-0.5em}
\begin{solcard}{A1}{Remove dependency $B(x)$. Take measurements at points on axis $Ox$ over the entire range $15~\text{mm} \leq x \leq 100~\text{mm}$. To secure the magnet on the table in a convenient position, you can use supports from the stand and adhesive mass. Conversion of voltage to field induction is not required at this point.}{2.00}

We fix the magnet in the stands and fix them on the table using adhesive mass. Let's place the sensor on the table at the height of the magnet at some distance from it. By rotating the magnet around its axis, we will achieve maximum voltmeter readings. This will ensure the correct orientation of the $Ox$ axis and its passage through the sensor. Let's remove the dependence $U(x)$.

\begin{figure}[H]
  \centering
  \includegraphics[width=0.72\textwidth]{E3_S_EN_assets/fig1.jpg}
\end{figure}
\end{solcard}

\begin{solcard}{A2}{Write down the exact formula $B(x)$ for a point magnetic dipole.}{0.60}

The general formula $\vec B \left(\vec r \right) = \cfrac{\mu_0}{4\pi}\cfrac{3(\vec m \cdot \vec r)\vec r - \vec m r^2}{r^5}$ in the case of $\vec m \upuparrows \vec r$ simplifies to the form $B(x)=\cfrac{\mu_0}{4\pi}\cfrac{2m}{x^3}$.

\par\smallskip\noindent\textbf{Answer:} $B(x)=\cfrac{\mu_0}{4\pi}\cfrac{2m}{x^3}$
\end{solcard}

\begin{solcard}{A3}{Plot a plot of $(-\ln B)$ versus $\ln x$ for the entire range of measured values. At what characteristic $x_{crit}$ does the nature of the dependence change?}{1.00}

The points of the graph corresponding to $x \in [15; 30]~\text{mm}$ fit well on a straight line with a slope coefficient $n = 2.02 \approx 2$, points at $x\in [55;100]~\text{mm}$ - on a straight line with a slope coefficient $3.03 \approx 3$, as for a point dipole. The nature of the dependence changes in the vicinity of $x_{crit}\approx 35~\text{mm}$.

\par\smallskip\noindent\textbf{Answer:} $x_{crit}\approx 35~\text{mm}$
\begin{figure}[H]
  \centering
  \includegraphics[width=0.72\textwidth]{E3_S_EN_assets/fig2.jpg}
\end{figure}
\end{solcard}

\begin{solcard}{A4}{Using the graph constructed in step A3, determine the value of $n$.}{0.40}
\par\smallskip\noindent\textbf{Answer:} $n = 2.02 \approx 2$
\end{solcard}

\begin{solcard}{A5}{Construct a linearized graph in such coordinates as to determine the magnetization $M$ of the magnet (magnetic moment per unit volume) from the slope of the graph. Calculate the numerical value of the quantity $M$.}{1.00}

Using the sensitivity of the sensor $k = 31.25~\text{V}/\text{T}$, we calculate the values of the field $B$. Let's construct a graph for large $x\geq60~\text{mm}$, where the previous graph shows that the points fit well on the second straight line. The slope coefficient $\beta$ of the graph $B$ from $1/x^3$, based on the theoretical formula, is $\beta= \cfrac{\mu_0 m}{2\pi} = \cfrac{\mu_0 \pi R^2 L}{2\pi} = \cfrac{\mu_0 MR^2 L}{2}$. According to the schedule, we determine $\beta=4.91 \cdot 10^{-7}$, from where $M=\cfrac{2\beta}{\mu_0 R^2 L}=7.81\cdot10^5~\text{A}/\text{m}$.

\par\smallskip\noindent\textbf{Answer:} $M=\cfrac{2\beta}{\mu_0 R^2 L}=7.81\cdot10^5~\text{A}/\text{m}$
\begin{figure}[H]
  \centering
  \includegraphics[width=0.72\textwidth]{E3_S_EN_assets/fig3.jpg}
\end{figure}
\end{solcard}

\begin{solcard}{B1}{Measure the dependence $h(d)$. Indicate the measurement errors, taking into account that the manufacturing accuracy of each lead is $\Delta d = \pm 0.05~\text{mm}$. Describe the type of dependence $h(d)$. There is no need to create a schedule.}{1.00}

Measurements of $h$ for all four available $d$ give the same result, determined very roughly due to the large relative error: $h(d)=(2.5 \pm 0.5)~\text{mm}$. Within the limits of measurement accuracy, there is no dependence on $d$. To illustrate that $h$ does not depend on $d$, we can place several different leads one after another in the “trough” and make sure that they line up in one line.

\par\smallskip\noindent\textbf{Answer:} $h(d)=(2.5 \pm 0.5)~\text{mm} \approx const$
\begin{figure}[H]
  \centering
  \includegraphics[width=0.72\textwidth]{E3_S_EN_assets/fig4.jpg}
\end{figure}
\end{solcard}

\begin{solcard}{B2}{On your answer sheet, note what type of material graphite is according to the observed effects: diamagnetic, paramagnetic, or ferromagnetic. Explain your choice.}{0.50}

The stylus is magnetized in the field of magnets and repels from them, trying to move to a position with a lower value of the magnet field $B$. This means that graphite is diamagnetic.

\par\smallskip\noindent\textbf{Answer:} Graphite is diamagnetic.
\end{solcard}

\begin{solcard}{B3}{Express the vertical force $F(y)$ acting on the magnetized stylus from the magnetic field in terms of the function $B(y)$ and the quantities $\chi$, $\mu_0$, $r$ and $l$. An expression containing other quantities will not be counted.}{1.00}

$F = -\cfrac{dU}{dy} = -\cfrac{d \left( -\cfrac{1}{2}(\vec m \cdot \vec B) \right)}{dy} = \cfrac{1}{2} \cfrac{2}{\mu_0} \cfrac{\chi}{\chi + 2} \pi r^2 l \cdot \cfrac{d}{dy} B^2 (y)$

\par\smallskip\noindent\textbf{Answer:} $F(y) = \cfrac{\pi r^2l}{\mu_0} \cfrac{\chi}{\chi+2} \cdot \cfrac{d}{dy} B^2 (y)$
\end{solcard}

\begin{solcard}{B4}{Obtain the final expression for $F(u)$ in terms of the quantities $\chi$, $\mu_0$, $M$, $r$, $l$, $R$, $a$, $u$. An expression containing other quantities will not be counted.}{1.00}

We substitute into the resulting formula $F(y)$ the explicit form of the function $ B(y) = \cfrac{\mu_0 M R^2}{a^2} \cdot \cfrac{1-u^2}{\left( 1+u^2 \right) ^2}$, where $u=\cfrac{y}{a}$. To take the derivative, it is convenient to write $ \cfrac{d}{dy} B^2 (y) = \cfrac{du}{dy} \cdot \cfrac{d}{du} B^2 (u) = \cfrac{1}{R} \cdot \cfrac{d}{du} B^2 (u) = \cfrac{1}{R} \cdot \cfrac{4u(1-u^2)(3-u^2)}{(1+u^2)^5}$. We arrive at the final answer:

\par\smallskip\noindent\textbf{Answer:} $F(u)=-\cfrac{\chi}{\chi+2} \cfrac{\pi r^2l \mu_0 M^2R^4}{a^5} \cfrac{4u(1-u^2)(3-u^2)}{(1+u^2)^5}$
\end{solcard}

\begin{solcard}{B5}{Write down how $u$ is expressed in terms of the measured value $h$ and what numerical value $u$ you will use in your calculations.}{0.50}
\par\smallskip\noindent\textbf{Answer:} $u=\cfrac{R-h}{a}=0.435$
\end{solcard}

\begin{solcard}{B6}{Calculate numerically the magnetic susceptibility of graphite $\chi$ based on the measurements from Part B1. Estimate the error of the result using the bounds method.}{1.00}

In the equilibrium position of the stylus (with measured $h$), we can write that the sum of forces along the vertical axis is zero: $$F(u)-m_0g = -\cfrac{\chi}{\chi+2} \cfrac{\pi r^2l \mu_0 M^2R^4}{a^5} \cfrac{4u(1-u^2)(3-u^2)}{(1+u^2)^5} - m_0 g = 0,$$ where $m_0 = \rho \pi r^2 l$. It can be seen, that $r$ cancels out, therefore theoretically confirmed, that position equilibrium $h$, implicitly entering in this formula, not depend from diameter of the lead $d$. Substitute $u=0.435$, $M=7.81\cdot10^5~\text{A}/\text{m}$ and lead to in condition quantity, obtain negative numerical value $\chi=-2.6\cdot 10^{-4} < 0$, that correspond a diamagnetic material. For estimate error one can use method boundary and substitute $h_{min} = 2.0~\text{mm}$ and $h_{max}=3.0~\text{mm}$. Taking into account value $\chi$, be obtained for such $h$, we can write:

\par\smallskip\noindent\textbf{Answer:} $\chi=(-2.6\pm0.7)\cdot 10^{-4}$
\end{solcard}

\begin{solcard}{C1}{Measure the required quantities and obtain the dependence $z_0(\alpha)$ for at least 4 different values $\alpha$. There is no need to create a schedule at this point. If you discover any peculiarity in the behavior of the stylus, write down the corresponding $\alpha_{crit}$.}{1.00}

We will tilt the table by tightening one of the support screws. By measuring the height at which the edge of the table closest to this screw is, we find the corresponding inclination of the table $\alpha=\cfrac{|h_{\alpha}-h_{\alpha=0}|}{\approx 20~\text{cm}}$. As the table tilts, the equilibrium position of the stylus gradually shifts to the support located lower on the slope. At $\alpha_{crit}\approx0.0175~\text{rad}$ the equilibrium position disappears, the stylus stops oscillating and one of its ends rests against the support of the magnet stand.

\par\smallskip\noindent\textbf{Answer:} $\alpha_{crit}\approx0.0175~\text{rad}$
\begin{figure}[H]
  \centering
  \includegraphics[width=0.72\textwidth]{E3_S_EN_assets/fig5.jpg}
\end{figure}
\end{solcard}

\begin{solcard}{C2}{Estimate the value of $\Delta U$ numerically and present your reasoning.}{1.00}

With a slope $\alpha_{crit}$, the difference in specific potential energies in the magnetic field between points $z=0$ and, for example, $z_0(\alpha_{crit})$, is compensated by the difference in specific potential energies in the gravitational field. As a result, the “dip” of potential disappears, the equilibrium position no longer exists, and the stylus moves to the position of the minimum total potential - near the support. Thus, we can estimate $\Delta U \sim g\alpha_{crit}z_0(\alpha_{crit}) = 1.4~\text{mJ}/\text{kg}$, where $z_0(\alpha_{crit}$ is the last measured equilibrium position before the stylus flies away to the support. A more accurate estimate is through the value $z_{max}=14~\text{mm}$ calculated in point C4 - the position of the “hump” of the potential: $\Delta U \sim g\alpha_{crit}z_{max} = 2.4~\text{mJ}/\text{kg}$.

\par\smallskip\noindent\textbf{Answer:} $\Delta U =2.4~\text{mJ}/\text{kg}$
\end{solcard}

\begin{solcard}{C3}{In a strictly horizontal position of the table (check by the equilibrium position of the stylus), measure the dependence $T(l)$ of the oscillation period on the length of the stylus in the maximum possible range of lengths $l$. Get at least 10 experimental points $T(l)$. There is no need to create a schedule at this point. Write down the critical lengths $l_{min}$ and $l_{max}$, at which the nature of the observed processes changes, and describe what is happening.}{1.50}

To increase the accuracy of determining $T$ we will measure the value $NT$ - the total time $N$ of oscillations. It is wise to choose $N\sim10...20$ because with $N>20$ the vibrations are almost imperceptible due to viscosity, and the final vibrations are almost impossible to measure.

\par\smallskip\noindent\textbf{Answer:} $l_{min} =8~\text{mm}$ – the stylus hangs crookedly or stands across the magnets. $l_{max} =28~\text{mm}$ – the stylus jumps out of its equilibrium position and crashes into the support of the stand.
\begin{figure}[H]
  \centering
  \includegraphics[width=0.72\textwidth]{E3_S_EN_assets/fig6.jpg}
\end{figure}
\end{solcard}

\begin{solcard}{C4}{Estimate the value of $z_{max}$ numerically and present your reasoning.}{0.50}

For lengths greater than $l_{max} =28~\text{mm}$, the position of stable equilibrium disappears. This happens because the ends of the stylus find themselves in areas of a very sharply decreasing magnetic field potential, and when the stylus is displaced, it is energetically favorable for the stylus to completely move away from the equilibrium position towards one of the supports.

\par\smallskip\noindent\textbf{Answer:} $z_{max} =\cfrac{l_{max}}{2}=14~\text{mm}$
\end{solcard}

\begin{solcard}{C5}{Based on the measurements of $z_0(\alpha)$ made in part C1, construct a linearized graph and determine the numerical value of the quantity $c_2$.}{1.00}

The total potential energy due to the magnetic and gravitational fields is equal to $U_\text{tot}(z) \approx c_0+c_2z^2+g\alpha z$. Correction $c_4z^4$ can be neglected here, as discussed earlier. This neglect is also justified by the linearity of the resulting linearized graph. Equilibrium position – minimum total potential energy: $\cfrac{dU_\text{tot}}{dz} \Bigg|_{z_0} \approx 2c_2z_0 + g \alpha = 0$. For linearization, we can plot $g \alpha$ against $2z_0$, then $c_2$ will be its slope coefficient.

\par\smallskip\noindent\textbf{Answer:} $c_2 =9.3~1/\text{s}^2$
\begin{figure}[H]
  \centering
  \includegraphics[width=0.72\textwidth]{E3_S_EN_assets/fig7.jpg}
\end{figure}
\end{solcard}

\begin{solcard}{C6}{Write down the law of conservation of energy for vibrations of a stylus of length $l$ in potential $U(z) = c_0 + c_2 z^2 + c_4 z^4$ and obtain a theoretical expression for the period $T(l)$ of small vibrations. The middle of the stylus is located at the point with coordinate $z=0$.}{1.00}

Let the stylus be displaced relative to the equilibrium position by a small distance $\Delta z$. The change in its potential energy in a magnetic field compared to the equilibrium position is equal to $$\Delta W = \rho \pi r^2 \Delta z \left(U\left(-\cfrac{l}{2}+\cfrac{\Delta z}{2}\right) - U\left(\cfrac{l}{2}+\cfrac{\Delta z}{2}\right)  \right) =\rho \pi r^2 \Delta z \cdot \left( \cfrac{dU}{dz} \Bigg|_{z=l/2} \cdot \Delta z\right)=$$ $$= \rho \pi r^2 (\Delta z)^2 \cdot (2c_2z+4c_4z^3)\Big|_{z=l/2}=\cfrac{\rho \pi r^2 l (\Delta z)^2}{2}\cdot (2c_2+c_4l^2) $$. Kinetic energy of the lead at the same time equal $K=\cfrac{\rho \pi r^2 l \cdot \dot{(\Delta z)}^2}{2}$. From energy conservation follow $dW+dK=0$: $$\cfrac{\rho \pi r^2 l\cdot 2(\Delta z) \dot{(\Delta z)}}{2}\cdot (2c_2+c_4l^2) + \cfrac{\rho \pi r^2 l \cdot 2\dot{(\Delta z)}\ddot{(\Delta z)}}{2}=0.$$ This expression one can cancel on $\rho \pi r^2 l\cdot \dot{(\Delta z)}$, remain: $$\ddot{(\Delta z)}=-(2c_2+c_4l^2)  \cdot (\Delta z).$$ Hence the oscillation period:

\par\smallskip\noindent\textbf{Answer:} $T(l)=\cfrac{2\pi}{\sqrt{2c_2+c_4l^2}}$
\end{solcard}

\begin{solcard}{C7}{Based on the measurements of $T(l)$ made in part C3, construct a linearized graph and determine the numerical values of the quantities $c_2$ and $c_4$.}{1.00}

To linearize the dependence at large lengths $22~\text{mm} \leq l \leq28~\text{mm}$, where the correction $c_4 z^4$ appears and the period significantly depends on $l$, we construct the dependence $\left( \cfrac{2\pi}{T} \right) ^2$ on $l^2$. The slope of the graph will be $c_4$ and the intercept will be $2c_2$.

\par\smallskip\noindent\textbf{Answer:} $c_2 = 23~1/\text{s}^2$, $c_4 = -4.51\cdot 10^4~\cfrac{1}{\text{m}^2 \text{s}^2}$
\begin{figure}[H]
  \centering
  \includegraphics[width=0.72\textwidth]{E3_S_EN_assets/fig8.jpg}
\end{figure}
\end{solcard}

\begin{solcard}{D1}{Carry out the necessary measurements of damped oscillations, construct a linearized graph and calculate numerically the value of $\gamma$.}{3.00}

If the resistance force is determined by the expression $F_\text{drag}/m = -\gamma \dot z$, then the vibration equation is written in the form $\ddot z + \gamma \dot z + \omega_0^2 z = 0$. The solution to this equation is harmonic oscillations with an amplitude that decays with time as $A(t)=A_0 e^{-\cfrac{\gamma}{2}t}$. The amplitude depends on the vibration number: $A(N)=A_0 e^{-\cfrac{\gamma}{2}NT}$. For linearization, you can plot the dependence of $\ln\left(\cfrac{A_0}{A(N)}\right)$ on $\cfrac{NT}{2}$, its slope coefficient will is equal to to $\gamma$. Let's measure the period of oscillations by conducting several series of 10...15 oscillations: $T=1.598~\text{s}$. Let's start damped oscillations by deflecting the stylus from the equilibrium position. We will record all successive amplitude values, observing them from a piece of graph paper inserted into the gap between the magnets. Let's carry out several such “runs”, plotting the corresponding series of points on a linearized graph. If different numbers of oscillations had approximately the same amplitudes, then on the graph we will mark only the average of the numbers. By averaging the slopes of the linear graphs obtained for different series of measurements, we obtain the average value $\gamma \approx 0.19~1/\text{s}$.

\par\smallskip\noindent\textbf{Answer:} $\gamma=0.19~1/\text{s}$
\begin{figure}[H]
  \centering
  \includegraphics[width=0.72\textwidth]{E3_S_EN_assets/fig9.jpg}
\end{figure}
\begin{figure}[H]
  \centering
  \includegraphics[width=0.72\textwidth]{E3_S_EN_assets/fig10.jpg}
\end{figure}
\end{solcard}

\end{document}